% answers_short.tex —— 2007 年·普通高等学校招生考试·陕西卷·理科 · 简短答案速查
% 由 tools/make_short_answers.py 自 answers.tex 的 \gksolhead/\gkanswer 提取，
% 并与 2007/2007_数学真题_陕西理_答案要点.md 逐题比对后生成，勿手改（改源 MD 后重跑）。

\begin{gknotice}
\textbf{使用说明}：本节只给最终答案便于快速核对，完整推导见解析版。答案为编写者独立推导，\textbf{非官方标准答案}。
\end{gknotice}

\gkpart{选择题}
\begin{center}\begin{tabular}{ccccccccccccc}
\toprule
题号 & 1 & 2 & 3 & 4 & 5 & 6 & 7 & 8 & 9 & 10 & 11 & 12 \\
\midrule
答案 & D & C & A & B & B & C & D & D & B & D & A & C \\
\bottomrule
\end{tabular}\end{center}

\gkpart{填空题}
\begin{itemize}[leftmargin=2.6em]
  \item[\textbf{13.}] $\dfrac13$。
  \item[\textbf{14.}] $8$。
  \item[\textbf{15.}] $6$。
  \item[\textbf{16.}] $210$。
\end{itemize}

\gkpart{解答题}
\begin{itemize}[leftmargin=2.6em]
  \item[\textbf{17.}] （1）$m=1$；（2）最小值 $1-\sqrt2$，此时 $x=-\dfrac{3\pi}{8}+k\pi$（$k\in\mathbf Z$）。
  \item[\textbf{18.}] （1）$\dfrac{101}{125}$；（2）分布列见解析，$E\xi=\dfrac{57}{25}$。
  \item[\textbf{19.}] （1）证明见解析；（2）$\arccos\dfrac{9}{\sqrt{93}}$。
  \item[\textbf{20.}] （1）$0<a<4$；（2）见解析。
  \item[\textbf{21.}] （1）$\dfrac{x^2}{3}+y^2=1$；（2）$\dfrac{\sqrt3}{2}$。
  \item[\textbf{22.}] （1）$a_k=k$；（2）$\dfrac1n$。
\end{itemize}
